% Shared scientific formulation for the manuscript and methodology companion.
\subsubsection{Relation to Yoon's MOCU}
For a specified operational loss $L(\theta,u)$, the model-specific and
intrinsically Bayesian robust (IBR) actions are, respectively,
\begin{align}
u_\theta^*&\in\arg\min_{u\in\mathcal U}L(\theta,u),\\
u_b^*&\in\arg\min_{u\in\mathcal U}\mathbb E_b[L(\theta,u)].
\label{eq:bayesian-control}
\end{align}
Yoon et al.'s construction compares these actions using the same loss
\cite{yoon2013mocu}; the generalized formulation also treats noisy
experiments~\cite{boluki2018generalized}:
\begin{equation}
\operatorname{MOCU}(b)=\mathbb E_b[L(\theta,u_b^*)-L(\theta,u_\theta^*)].
\label{eq:general_mocu}
\end{equation}
Bayesian robustness means minimizing expected cost, not necessarily taking a
worst-case maximum. The loss below is our power-grid specialization; neither
its coefficient nor the 95\% level is prescribed by Yoon et al.

\subsubsection{Physical control requirement and finite loss}
Let $\mathcal U=[0,u_{\max}]$ be the continuous terminal-control range. Define
$U(\theta)$ as the minimum magnitude satisfying the prescribed nadir and
ROCOF limits for every $c\in\mathcal C$. The scalar requirement model assumes
that safety is monotone over this range and that a safe magnitude exists.
These assumptions must be checked for each simulator and contingency set;
otherwise a scalar threshold is insufficient and the full simulated
action-dependent loss must be used. Equivalently, under these assumptions,
\begin{equation}
\mathcal U_{\mathrm{safe}}(\theta)=[U(\theta),u_{\max}].
\label{eq:safe_control_set}
\end{equation}
The current IEEE9 scenario is the configured single contingency; a general
set $\mathcal C$ does not establish N--1 security without evaluating that set.

For $\lambda>1$ and $(z)_+=\max(z,0)$, choose
\begin{equation}
L(\theta,u)=u+\lambda[U(\theta)-u]_+,
\qquad \lambda=20.
\label{eq:operational_cost}
\end{equation}
This finite loss penalizes supplementary injection and unmet required
injection. It is a power-equivalent operational surrogate, not an empirically
calibrated monetary damage model or a direct measure of violation severity.
The coefficient is prespecified to encode an asymmetric preference against
under-control. A sensitivity study should vary this preference while changing
the decision quantile consistently. No additional binary violation penalty
or control margin is included in the primary formulation.

Since the loss decreases with slope $1-\lambda$ below $U$ and increases with
slope $1$ above $U$, its model-specific minimizer is
\begin{equation}
u_\theta^*=u_{\mathrm{optimal}}(\theta)=U(\theta),
\quad L(\theta,u_\theta^*)=U(\theta).
\label{eq:model_specific_control}
\end{equation}

\subsubsection{Derivation of the 95\% decision}
Write $F_b(v)=\Pr_b(U\le v)$ and $\alpha=1/\lambda=0.05$.
The left and right derivatives of posterior expected loss at $u$ are
$1-\lambda+\lambda F_b(u^-)$ and
$1-\lambda+\lambda F_b(u)$, respectively. Convexity therefore gives the
optimality condition
\begin{equation}
F_b(u^- )\le 1-\alpha\le F_b(u).
\label{eq:quantile_optimality}
\end{equation}
Choosing the smallest minimizer yields
\begin{equation}
u_b^*=Q_{0.95}(U\mid b)
=\inf\{v:F_b(v)\ge0.95\}.
\label{eq:u_ctrl_support}
\end{equation}
Thus the quantile is derived from the finite loss, rather than appended to
an unrelated MOCU definition. Equivalently, realized regret is
\begin{equation}
\operatorname{OCU}(\theta,b)
=(u_b^*-U)_++19(U-u_b^*)_+\ge0.
\label{eq:ocu}
\end{equation}
Over-control has unit marginal regret and under-control has marginal regret
19; the coefficient 20 in the total loss includes the avoided injection cost.
The corresponding belief MOCU is
\begin{equation}
\operatorname{MOCU}(b)=\mathbb E_b\!\left[
u_b^*+20(U-u_b^*)_+-U\right].
\label{eq:belief_mocu}
\end{equation}
It vanishes when the required control is known exactly, even if other model
parameters remain uncertain. Subtracting $\mathbb E_b[U]$ without including
the shortfall term is not this finite-loss MOCU.

\subsubsection{Particle bank and implemented terminal action}
Let $U_n$ denote the stored control requirement for particle $\theta^{(n)}$,
with normalized posterior weights $w_n$. The empirical CDF and decision are
\begin{align}
F_N(v)&=\sum_n w_n\mathbf1\{U_n\le v\},\\
q_N&=\inf\{v:F_N(v)\ge0.95\},\\
u_{\mathrm{ctrl}}(b)&=\operatorname{snap}_{\mathcal G}^{\uparrow}(q_N).
\label{eq:u_ctrl}
\end{align}
Here $\mathcal G$ is the configured control grid. The aligned implementation
requires every $U_n\in\mathcal G$ and zero margin, so the empirical quantile
already lies on the grid and snapping does not change it. It is consequently
an exact IBR action for this finite particle-bank loss. If requirements lie
off-grid, rounding a continuous quantile upward is not generally the
unconstrained grid Bayes optimum; one must minimize expected loss over the
grid explicitly or retain continuous controls. The implementation checks
this alignment before running the primary experiment.

The bank-level objective used for training and posterior evaluation is
\begin{equation}
\widehat{\operatorname{MOCU}}_N(b)
=\sum_n w_n\{u_{\mathrm{ctrl}}+20(U_n-u_{\mathrm{ctrl}})_+-U_n\}.
\label{eq:particle_mocu}
\end{equation}
All methods use the same bank, loss, quantile, and control range. Probe and
control bank rows must match in particle values and ordering, physics,
scenario, units, and simulator provenance. A dense control-grid bank is a
numerical approximation to the continuous physical requirements; its error
must be distinguished from posterior uncertainty.

\subsubsection{Sequential design and held-out validation}
The ideal design objective and its implemented approximation are
\begin{align}
\pi^*&\in\arg\min_\pi\mathbb E_{H_T\mid\pi}
[\operatorname{MOCU}(b_T)],\\
\widehat\pi&\approx\arg\min_\pi\mathbb E_{H_T\mid\pi}
[\widehat{\operatorname{MOCU}}_N(b_T)].
\label{eq:mocu_policy}
\end{align}
Under an exact posterior and the same generative model, the tower property
equates expected posterior MOCU to expected realized regret. Also,
$\mathbb E_{H_T}\mathbb E_{b_T}[U]=\mathbb E_{p_0}[U]$ is independent of the
probe policy: minimizing expected MOCU is equivalent to minimizing expected
total operational loss, not merely the control magnitude. Optimal expected
MOCU cannot increase when additional observations may be ignored; this does
not imply monotonic improvement for every realized history or trained policy.

For an independent test system $\theta^\star$, evaluate the applied action
against a separately computed continuous physical oracle $U^*(\theta^\star)$:
\begin{equation}
R^*=u_{\mathrm{ctrl}}+20[U^*(\theta^\star)-u_{\mathrm{ctrl}}]_+
-U^*(\theta^\star).
\label{eq:realized_regret}
\end{equation}
Mean held-out $R^*$ is the primary empirical endpoint; posterior bank MOCU
is reported separately. Their difference can reflect finite support,
control-grid approximation, and model or posterior miscalibration. Regret
against a continuous oracle includes control discretization error and should
not be described as purely uncertainty cost for a restricted grid action
class. A strictly grid-relative MOCU instead subtracts
$\min_{v\in\mathcal G}L(\theta,v)$ for each model.

Under the threshold model and exact posterior, the quantile ensures posterior
coverage $\Pr_b(U\le u_b^*)\ge0.95$. It does not protect every posterior
model and is not a guaranteed empirical safety rate. Report true-system
nadir and ROCOF violations, safety counts and confidence intervals, control
magnitude, and regret separately. Any 95\% safety eligibility rule must be
prespecified; retain results for ineligible methods and distinguish a point
estimate passing 95\% from statistical evidence for that claim. The finite
loss is scientifically a MOCU specialization, while its engineering adequacy
depends on the stated loss preference, physical validation, and calibration.
